\section*{अध्याय \ref{ch:g9:microbes-and-infection} --- सूक्ष्मजीव और संक्रमण}

\begin{solution}{exo:g9:microbes-and-infection:1}
\omterm{def:g9:microbes-and-infection:sorted}{जीवाणु} (बिना \omterm{def:g6:cell-unit-of-life:parts}{केंद्रक} वाली अकेली \omterm{def:g5:first-look-at-cells:cell}{कोशिकाएँ}), ख़मीर जैसे \omterm{def:g6:cell-unit-of-life:unicellular}{एककोशिकीय} \omterm{def:g6:decomposers-and-soil:decomposers}{कवक} और
तालाब का बाक़ी \omterm{def:g6:cell-unit-of-life:unicellular}{एककोशिकीय} जीवन --- और \omterm{def:g9:microbes-and-infection:sorted}{विषाणु}, जो क़िस्म में ही अलग हैं: न
कोई \omterm{def:g5:first-look-at-cells:cell}{कोशिका}, न अपना कोई जीवन, बस बँधी हुई आनुवंशिक इबारत, जो किसी जीवित
\omterm{def:g5:first-look-at-cells:cell}{कोशिका} की मशीनरी उधार ले लेती है।
\end{solution}

\begin{solution}{exo:g9:microbes-and-infection:2}
वे खरबों बेज़रर \omterm{def:g9:microbes-and-infection:sorted}{जीवाणु}, जो \omterm{def:g1:the-five-senses:sense}{त्वचा} पर और आँत में बसे हैं। सेवाएँ: \omterm{def:g4:journey-of-food:digestion}{पाचन} में
मदद, और ज़मीन घेरना --- यानी \omterm{prop:g9:microbes-and-infection:mostly}{रोगजनकों} को उतरने की जगहों से हटाए रखना।
\end{solution}

\begin{solution}{exo:g9:microbes-and-infection:3}
संचरण (किसी खाँसी की बूँदें); प्रवेश (कोई कटाव, हवा-मार्ग); बढ़त
(घंटे-दर-घंटे दोगुने होते बँटवारे, या हड़पी हुई \omterm{def:g5:first-look-at-cells:cell}{कोशिकाओं} की नक़ल); और
बीमारी (पकी हुई उँगली की टीस, ज़ुकाम का कच्चा गला)।
\end{solution}

\begin{solution}{exo:g9:microbes-and-infection:4}
कटाव का \omterm{def:g9:microbes-and-infection:sorted}{जीवाणु} गरम गीले ऊतक में \omterm{def:g5:first-look-at-cells:cell}{कोशिकाओं} के बीच बढ़ता है और अपना
आस-पड़ोस ज़हरीला करता है। और ज़ुकाम का \omterm{def:g9:microbes-and-infection:sorted}{विषाणु} हवा-मार्ग की परत की
\omterm{def:g5:first-look-at-cells:cell}{कोशिकाओं} के भीतर बढ़ता है, और हर हड़पी हुई \omterm{def:g5:first-look-at-cells:cell}{कोशिका} नक़लों की एक भीड़ छोड़
देती है।
\end{solution}

\begin{solution}{exo:g9:microbes-and-infection:5}
\omterm{def:g1:the-five-senses:sense}{त्वचा}: प्रवेश पर वह दीवार। परतों की बुहारियाँ और स्नान (हवा-मार्ग की
सफ़ाई, आँसू, \omterm{def:g4:journey-of-food:tract}{आमाशय} का अम्ल): प्रवेश और शुरुआती बढ़त। \omterm{prop:g9:microbes-and-infection:mostly}{बसे हुए
सूक्ष्मजीव}: प्रवेश --- ज़मीन घेरी हुई। और \omterm{def:g1:healthy-habits:hygiene}{स्वच्छता} का साबुन, पकाना तथा
ठंडक: संचरण और बढ़त।
\end{solution}

\begin{solution}{exo:g9:microbes-and-infection:6}
संचरण --- यानी वह मुसाफ़िर, जिसे डॉक्टरों के हाथ शव-परीक्षा के कमरे से
प्रसव-कक्ष तक ढो रहे थे। और पास्तूर ने ख़ुद उन मुसाफ़िरों को दिखा दिया:
खट्टा होने, बीमारी और \omterm{def:g9:microbes-and-infection:stages}{संक्रमण} के कारक के रूप में \omterm{def:g9:microbes-and-infection:sorted}{सूक्ष्मजीव} --- और
उन्हें रोकने वाली गरमी।
\end{solution}

\begin{solution}{exo:g9:microbes-and-infection:7}
प्रतिजैविक \omterm{def:g9:microbes-and-infection:sorted}{जीवाणुओं} को भूखा मारते और तोड़ते हैं --- यानी उन आज़ाद
\omterm{def:g5:first-look-at-cells:cell}{कोशिकाओं} को, जिनकी अपनी मशीनरी में ज़हर घोला जा सकता है। जबकि ज़ुकाम के
\omterm{def:g9:microbes-and-infection:sorted}{विषाणु} की अपनी कोई मशीनरी है ही नहीं: वह \emph{तुम्हारी} पर चलता है, और
उसकी कार्य-प्रणाली पर साधी गई दवा तुम्हारी \omterm{def:g5:first-look-at-cells:cell}{कोशिकाओं} पर ही निशाना साधती
है --- इसीलिए \omterm{def:g9:microbes-and-infection:sorted}{जीवाणु} की दवा \omterm{def:g9:microbes-and-infection:sorted}{विषाणु} को अनछुआ छोड़ देती है।
\end{solution}

\begin{solution}{exo:g9:microbes-and-infection:8}
खाने से पहले हाथ धोना: हाथ-से-मुँह वाला संचरण रोकता है। शौच के बाद धोना:
आँत वाले रास्ते का निकास-से-प्रवेश वाला फंदा रोकता है। खाँसी को छींक की
तरह ढककर लेना (और \omterm{def:g1:my-body:joint}{कोहनी} में लेना): स्रोत पर ही बूँदों का संचरण रोकता
है। (और दाँतों तथा भोजन के नियम मुँह की बढ़त वाली जगहें रोकते हैं।)
\end{solution}

\begin{solution}{exo:g9:microbes-and-infection:9}
\omterm{prop:g9:microbes-and-infection:mostly}{रोगजनक}: मुर्ग़ पर बैठे भोजन-विषाक्तता वाले \omterm{def:g9:microbes-and-infection:sorted}{जीवाणु}। संचरण: ख़ुद वह पकवान।
प्रवेश: आँत। किससे बढ़ा: गाड़ी के तीन गरम घंटे --- यानी परिस्थिति-नियम,
जो हर आधे घंटे में दोगुना होने को खिलाता रहा। सबसे सस्ती रोक: ठंडक की
शृंखला --- कोई ठंडा डिब्बा (या उसे ताज़ा ही खा लेना); और उसके बाद
सबसे सस्ती, पहले से अच्छी तरह पकाना।
\end{solution}

\begin{solution}{exo:g9:microbes-and-infection:10}
क्योंकि हाथ अलग-अलग ख़ानों के बीच हरकारे हैं: ज़ेमेलवाइस के डॉक्टर
रोज़मर्रा के किसी भी पैमाने पर साफ़ थे, फिर भी वार्ड का क़ातिल लाश से
माँ तक ढो रहे थे। जो धुलाई मायने रखती है वह \emph{सरहद पर} वाली है ---
कच्चे मुर्ग़ से सलाद तक, शौचालय से मेज़ तक --- यानी वहाँ, जहाँ वह रास्ता
वरना जुड़ जाता।
\end{solution}

\begin{solution}{exo:g9:microbes-and-infection:11}
वह प्रतिजैविक बढ़त पर हमला करता है --- यानी बँटते हुए \omterm{def:g9:microbes-and-infection:sorted}{जीवाणुओं} को तोड़ता
है। और अधूरे ढंग से इस्तेमाल हो तो वह आधे रास्ते तक चली एक चयन-रोक बन
जाता है: कमज़ोर मर जाते हैं, और रोधी अल्पसंख्या उस वार्ड की वारिस बन
जाती है --- यानी वही \omterm{prop:g9:microbes-and-infection:mostly}{रोगजनक} पाला जाता है जिससे अगले मरीज़ का सामना होगा।
\end{solution}

\begin{solution}{exo:g9:microbes-and-infection:12}
वह इबारत: आनुवंशिक अक्षरों का एक टुकड़ा, और वर्णमाला वही जो हर \omterm{def:g5:first-look-at-cells:cell}{कोशिका} की
है (\cref{rem:g9:genes-and-dna:universal})। वह छापाख़ाना: किसी जीवित
\omterm{def:g5:first-look-at-cells:cell}{कोशिका} की पढ़ने और बनाने वाली मशीनरी, जिसे भीतर घुसी वह इबारत अपनी ही
नक़लें छापने पर लगा देती है। जीवन का किनारा: अकेला \omterm{def:g9:microbes-and-infection:sorted}{विषाणु}
\cref{prop:g1:living-or-not:signs} की सूची में से कुछ भी नहीं करता ---
न खाना, न बढ़ना, न बँटना; और उधार पर वह जनन की एक नक़ल भर चलाता है। दवा
मुश्किल: क्योंकि ज़हर घोलने को कोई विषाणु-मशीनरी है ही नहीं --- सिर्फ़ वह
छापाख़ाना, और वह छापाख़ाना तुम हो।
\end{solution}

\begin{solution}{pb:g9:microbes-and-infection:1}
\textbf{1.} बूँदों वाला संचरण --- खाँसी और बोली से --- और साथ में
टेबल-टेनिस के बल्लों वाली साझी सतहों की मदद; और प्रवेश हवा-मार्गों से।

\textbf{2.} बढ़त का: हर मामले की हड़पी हुई \omterm{def:g5:first-look-at-cells:cell}{परत-कोशिकाएँ} ऐसी नक़लें छापती
हैं जो अगले मामलों को बो देती हैं --- यानी यह दोगुना होना उस \omterm{def:g9:microbes-and-infection:sorted}{विषाणु} की
उधार ली हुई मशीनरी है, जो जुड़ती चली जाती है।

\textbf{3.} बीमारी। और वे हड़पी हुई, नाकाम होती \omterm{def:g5:first-look-at-cells:cell}{परत-कोशिकाओं} के नुक़सान
से उठते हैं --- तथा कुछ हद तक ख़ुद हिफ़ाज़त की लड़ाई से, जैसा अगला अध्याय
ब्योरे में बताता है।

\textbf{4.} क्योंकि सभाएँ हर कक्षा की बूँद-पहुँच को घंटे भर के लिए जोड़
देती हैं: संचरण भीड़ के साथ बढ़ता है, और स्कूल के वे दो सबसे घने मेल ही
ख़ानों के बीच का रास्ता थे।

\textbf{5.} बीमार विद्यार्थी घर: स्रोत पर ही संचरण हटाया गया।
खिड़कियाँ: बूँदों की ख़ुराक पतली हुई --- संचरण। हाथ धोना: सतहों वाला
रास्ता कटा --- संचरण। सामान पोंछना: साझी चीज़ों वाला रास्ता कटा ---
संचरण। (और सब उसी सबसे सस्ते चरण पर साधे हुए।)

\textbf{6.} ``प्रतिजैविक \omterm{def:g9:microbes-and-infection:sorted}{जीवाणु} मारते हैं; और यह कारक एक \omterm{def:g9:microbes-and-infection:sorted}{विषाणु} है, जो
बच्चों की अपनी ही \omterm{def:g5:first-look-at-cells:cell}{कोशिकाओं} पर चल रहा है --- वह दवा उसे छू ही नहीं सकती,
और उसे फिर भी बिखेरने का नतीजा बस इतना होगा कि हर किसी के बसे हुए
\omterm{def:g9:microbes-and-infection:sorted}{जीवाणुओं} पर एक चयन-छननी चल जाए, और उस दिन के लिए रोध पल जाए जब
प्रतिजैविक सचमुच चाहिए होंगे।''

\textbf{7.} कारक: भोजन-विषाक्तता वाले \omterm{def:g9:microbes-and-infection:sorted}{जीवाणु}। संचरण: वह मिठाई। प्रवेश:
आँत। बढ़त: बिना फ़्रिज बीते घंटे --- यानी गरमाहट, नमी, भोजन और समय, सब
मिले हुए। रोक: फ़्रिज --- यानी \omterm{def:g6:exploring-our-environment:conditions}{परिस्थितियों} से इनकार; और उन दोनों फैलावों
में स्कूल के सिवा कुछ भी साझा नहीं।

\textbf{8.} फैलाव का नक़्शा संचरण को \omterm{def:g1:the-five-senses:sense}{आँखों} से दिखने लायक़ बना देता है:
कौन-कहाँ-बैठा उस रास्ते को खींच देता है --- क़तारें, साथी, सभाएँ --- और
वह रास्ता एक बार खिंच जाए तो दिखा देता है कि काटना कहाँ है। बेतरतीब इलाज
मामलों से लड़ता है; नक़्शा फैलाव से।

\textbf{9.} संचरण: सूचकांक मामले की बूँदें, पहले उसकी क़तार तथा साथी,
और फिर सभाओं से पूरा स्कूल। प्रवेश: हवा-मार्ग, कक्षा-दर-कक्षा। बढ़त:
पहले हफ़्ते भर दो-दो दिन में दोगुना होना। बीमारी और बुझना: बुख़ार चोटी पर
आते हैं, उपाय रास्ते काटते हैं, और पकड़ में आने लायक़ों का तालाब पतला
पड़ता है --- वक्र नीचे मुड़ जाता है।

\textbf{10.} उपायों की रोकों ने उन रास्तों को भूखा रखा --- और वह तालाब
ख़ुद भी हाज़िरी-रजिस्टर से छोटा था: जो विद्यार्थी पिछली सर्दी की चचेरी
नस्ल से प्रतिरक्षा लिए थे उन्हें जलाया ही नहीं जा सकता था, और कोई फैलाव
तभी मर जाता है जब पकड़ में आने लायक़ बहुत कम बचें: यानी अगले अध्याय की
याददाश्त, पहले से झलकती हुई।

\textbf{11.} ज़ेमेलवाइस का: हरकारों के रास्ते काटो --- घर पर अलगाव, धुले
हाथ, पोंछे हुए बल्ले। और पास्तूर का: उस मुसाफ़िर को पहचानो --- कारक का
नाम लो, उसे \omterm{def:g6:exploring-our-environment:conditions}{परिस्थितियाँ} देने से इनकार करो, और इलाज उसकी क़िस्म पर साधो,
रिवाज़ पर नहीं।

\textbf{12.} ``फैलाव सबसे सस्ते में संचरण पर रुकते हैं --- साबुन, हवा,
दूरी और एक बंद ठंडा डिब्बा उस हर दवा से बेहतर हैं जो बढ़त का इंतज़ार करती
है।''
\end{solution}
